\chapter{Физическая роль класса пятнадцать: архитектурное решение}
\label{ch:v6-spectral-transition-class15-physical-role-branch-decision-gate}

\section{Зачем требуется явное решение}

Класс
\begin{equation}
 \kappa_{15}=15\in KO_6
 \bigl(M_{105}(\mathbb C)_{\mathbb R}\bigr)
\end{equation}
строго построен в Томе V. Однако последовательность последних гейтов
разделила три возможных чтения одного и того же числа:
\begin{enumerate}
 \item классификационный коэффициент полного пакета поколения;
 \item число состояний одной связанной частицы;
 \item кратность уровней, пересекающих ноль в событии рождения.
\end{enumerate}
Первое чтение доказано. Второе и третье не следуют из $K$-класса
автоматически. Настоящий гейт запрещает дальнейшую подмену этих ролей.

\section{Строгий ledger входных результатов}

Архитектурное решение опирается на уже закрытые сертификаты:
\begin{enumerate}
 \item тёплицев и Real-класс пятнадцать существует;
 \item физическая минимальная опора раскладывается как $15=12+3$;
 \item оба слагаемых имеют собственные тёплицевы и KO6-подциклы;
 \item текущий родитель не колокализует ранги $12$ и $3$;
 \item хиггс-разрешённая опора распадается как $6+6+2+1$;
 \item один дублет не топологически принуждает локальный нуль Хиггса;
 \item сфалерон реализует $H=0$, но является неустойчивым седлом;
 \item физический сфалеронный поток равен $4$, а не $15$;
 \item формальный product даёт $15$ только после подстановки готового
 $q_0$, тогда как эквивариантный product снова даёт $4$.
\end{enumerate}

Теория индекса классифицирует устойчивые классы операторов и их
спаривания, но сама по себе не превращает класс в локализованную
устойчивую частицу или конкретную динамическую траекторию
\cite{class15-role-atiyah-singer1968,
class15-role-kasparov1981}. Аналогично, конечная спектральная тройка
фиксирует представления и разрешённые рёбра, а не автоматически
связывает все её состояния в один солитон
\cite{class15-role-krajewski1998}.

\section{Три варианта продолжения}

\subsection*{Вариант A: продолжать называть 15 кратностью рождения}

Этот вариант несовместим с физическим индексом $4$, разложением
$12+3$, отсутствием колокализации и отрицательной модой сфалерона. Он
требовал бы считать формальное равенство $1\cdot15=15$ динамическим
выводом, хотя число пятнадцать уже содержится во входном проекторе.

\subsection*{Вариант B: добавить переход полного конечного оператора}

Можно постулировать новый путь, на котором через ноль проходят также
правые слабые синглеты. Но такой путь не является стандартным
сфалероном. Для него потребуются новый самосопряжённый оператор,
действие, gauge-ковариантность, правила аномалий, KO6-совместимость,
граничные условия и нелинейное насыщение. Это допустимое направление
новой модели, но не результат текущей версии.

\subsection*{Вариант C: классификационный ledger и компонентная физика}

Класс пятнадцать сохраняется как глобальный паспорт полного
однопоколенного пакета:
\begin{equation}
 \kappa_{15}=\kappa_{12}+\kappa_3.
\end{equation}
Он кодирует суммарный коэффициент, Real-пару и вес $1/7$, но не
утверждает, что все пятнадцать направлений рождаются одним событием,
живут в одном пространственном ядре или образуют одну частицу.
Физические переходы должны проверяться на компонентных носителях и с
их фактическими представлениями.

\begin{theorem}[Роль класса пятнадцать в версии VI]
В текущей архитектуре класс пятнадцать является классификационным
ledger полного фермионного пакета одного поколения. Он не является
доказанной кратностью сфалеронного рождения, числом частиц, рангом
одного связанного солитона или числом физических нулевых пересечений.
Для такого усиления требуется явно новая динамическая модель.
\end{theorem}

\section{Выбранная ветвь}

В версии VI выбирается вариант C. Это решение сохраняет все строгие
положительные результаты:
\begin{itemize}
 \item класс $15$ и его Real/Toeplitz-факторизацию;
 \item компонентные классы $12$ и $3$;
 \item хиггс-разрешённые опоры $6+6+2+1$;
 \item стандартный сфалеронный поток $4$ как независимый физический
 контроль.
\end{itemize}
Одновременно снимаются недоказанные утверждения о единой частице и
едином событии рождения пятнадцати состояний.

Это не означает, что пятнадцать фермионных направлений не могут иметь
общего происхождения в более глубокой модели. Установлено только, что
такое происхождение не выведено текущими операторами Тома VI.

\begin{nogocriterion}
Любой последующий текст версии VI, использующий число пятнадцать как
физическую multiplicity рождения, обязан предъявить новый оператор и
пройти повторные тесты представлений, аномалий, Real/KO6, локализации и
устойчивости. Ссылка только на ранг $q_0$ или формальный внешний product
недостаточна.
\end{nogocriterion}

\section{Следующий проверяемый вопрос}

После отказа от единого события нужно определить, остаётся ли в текущем
родителе хотя бы одно параметрически чистое \emph{компонентное}
наблюдаемое рождения. Следующий гейт
\path{version6_spectral_transition_componentwise_creation_observable_gate}
должен сопоставить:
\begin{enumerate}
 \item независимые классы $\kappa_{12}$ и $\kappa_3$;
 \item физический слабый поток $(3,1)$;
 \item аномальные заряды $B,L,B-L$;
 \item локальные хиггсовские опоры $6+6+2+1$;
 \item наличие или отсутствие единой динамической амплитуды.
\end{enumerate}

Нулевая гипотеза: текущая архитектура фиксирует только допустимые
целочисленные и зарядовые selection rules, но не абсолютную вероятность,
массу или число рождаемых частиц.

Стоп-критерий: если ни один компонентный observable не следует без
нового масштаба, профиля или переходного оператора, программа рождения
материи версии VI фиксируется как незамкнутая на динамическом уровне.

\begin{protocol}
\begin{enumerate}
 \item класс $15$ как KO6/Toeplitz ledger: \textbf{сохранён};
 \item класс $15$ как неделимая опора: \textbf{опровергнут};
 \item класс $15$ как связанная частица: \textbf{не выведен};
 \item класс $15$ как сфалеронная multiplicity: \textbf{опровергнут};
 \item физический сфалеронный поток: \textbf{$4$};
 \item компонентные классы: \textbf{$12$ и $3$};
 \item формальный product $1\cdot15$: \textbf{классификационный};
 \item новая полная переходная динамика: \textbf{требует новой модели};
 \item выбранная ветвь версии VI:
 \textbf{ledger плюс компонентная физика};
 \item единое рождение наблюдаемого поколения: \textbf{не доказано};
 \item физическое замыкание: \textbf{не достигнуто}.
\end{enumerate}
\end{protocol}

Машинный ledger решения сохранён в
\path{s2t_v6_spectral_transition_class15_physical_role_branch_decision_gate_results.json}.

\begin{thebibliography}{9}
\bibitem{class15-role-atiyah-singer1968}
M.~F.~Atiyah, I.~M.~Singer,
\emph{The Index of Elliptic Operators: I},
Annals of Mathematics 87 (1968), 484--530.

\bibitem{class15-role-kasparov1981}
G.~G.~Kasparov,
\emph{The Operator $K$-Functor and Extensions of $C^*$-Algebras},
Mathematics of the USSR-Izvestiya 16 (1981), 513--572.

\bibitem{class15-role-krajewski1998}
T.~Krajewski,
\emph{Classification of Finite Spectral Triples},
Journal of Geometry and Physics 28 (1998), 1--30;
arXiv:hep-th/9701081.
\end{thebibliography}